\chapter{Structure complexe interne et géométrie projective des états à \(2\) niveaux}
\label{chap:geometria-proyectiva-bloch}
\index{racine interne de moins un}
\index{Paley!matrice de conférence}
\index{Witt}
\index{sphère de Bloch}
\index{Möbius!action}
\index{groupe icosaédrique binaire}

La géométrie de phase employée par la réalisation quantique possède une
généalogie algébrique antérieure à la notation complexe. Son point de départ est
l'incidence finie de Paley--Witt ; on en obtient un endomorphisme dont le
carré est moins l'identité. La réalification conserve cet opérateur comme
générateur de phase, et la projectivisation d'un espace à deux niveaux conduit
à la sphère de Bloch. Sur cette sphère, l'action spinorielle de
\(\mathrm{SU}(2)\), son expression par transformations de Möbius et le groupe
icosaédrique binaire forment des réalisations successives d'une même structure.

La chaîne causale développée dans ce chapitre est
\begin{equation}
 \boxed{
 \begin{gathered}
 C_P\longrightarrow A_W^2=-I_6,\\
 A_W^2=-I_6\longrightarrow
 \mathbb F_3^6\simeq\mathbb F_9^3
 \longrightarrow (E,J_E),\\
 (E,J_E)\longrightarrow
 \mathbb P(\mathcal H_J)\simeq\mathbb{CP}^{1}
 \longrightarrow S^2,\\
 S^2\longrightarrow
 \bigl\{\mathrm{SU}(2),\ \mathop{\text{M\"ob}}(S^2),\ 2A_5\bigr\}.
 \end{gathered}}
 \label{eq:cadena-paley-bloch-anexo}
\end{equation}
Chaque flèche déclare son domaine. La structure quadratique finie, la
réalification orthogonale et la géométrie projective restent liées, avec
leurs corps de scalaires explicitement séparés.

\section{L'incidence de Paley comme origine algébrique de la phase}

Soit \(C_P\in M_6(\mathbb Z)\) la matrice symétrique de conférence
\begin{equation}
 C_P=
 \begin{pmatrix}
 0& 1& 1& 1& 1& 1\\
 1& 0& 1&-1&-1& 1\\
 1& 1& 0& 1&-1&-1\\
 1&-1& 1& 0& 1&-1\\
 1&-1&-1& 1& 0& 1\\
 1& 1&-1&-1& 1& 0
 \end{pmatrix}.
 \label{eq:paley-conferencia-anexo}
\end{equation}
Ses lignes satisfont
\begin{equation}
 C_P C_P^{\mathsf T}=5I_6.
 \label{eq:paley-ortogonalidad-anexo}
\end{equation}
En réduisant modulo trois, nous obtenons
\begin{equation}
 A_W:=C_P\bmod 3\in M_6(\mathbb F_3).
 \label{eq:paley-reduccion-anexo}
\end{equation}

\begin{theorem}[Racine interne de moins l'identité]
\label{thm:raiz-interna-menos-uno-anexo}
L'endomorphisme \(J_W:v\mapsto vA_W\) de
\(V_W=\mathbb F_3^6\) satisfait
\begin{equation}
 J_W^2=-I_{V_W}.
 \label{eq:Jw-cuadrado-anexo}
\end{equation}
En conséquence, \(V_W\) admet une structure canonique de module sur
\(\mathbb F_3[X]/(X^2+1)\) une fois la matrice de Paley fixée.
\end{theorem}

\begin{proof}
La matrice \(C_P\) est symétrique. En réduisant
\eqref{eq:paley-ortogonalidad-anexo} modulo trois, on obtient
\[
 A_W^2=5I_6\bmod3=2I_6=-I_6.
\]
Ainsi, l'évaluation \(X\mapsto J_W\) annule l'idéal
\((X^2+1)\) et définit l'action du quotient indiqué.
\end{proof}

Le polynôme \(X^2+1\) est irréductible sur \(\mathbb F_3\), car ses
valeurs en \(0,1,2\) sont \(1,2,2\). Par conséquent,
\begin{equation}
 \mathbb F_9:=\mathbb F_3[\iota]/(\iota^2+1)
 \label{eq:F9-desde-paley-anexo}
\end{equation}
est un corps et \(J_W\) réalise la multiplication par \(\iota\).

\begin{corollary}[Descente de six coordonnées ternaires à trois coordonnées quadratiques]
\label{cor:F3seis-F9tres-anexo}
La règle
\begin{equation}
 (a+b\iota)v:=av+bJ_Wv,
 \qquad a,b\in\mathbb F_3,
 \label{eq:accion-F9-anexo}
\end{equation}
transforme \(V_W\) en un espace vectoriel de dimension trois sur
\(\mathbb F_9\). Dans les coordonnées fixées par \(A_W\),
\begin{equation}
 \mathbb F_3^6\simeq\mathbb F_9^3.
 \label{eq:F3seis-igual-F9tres-anexo}
\end{equation}
\end{corollary}

\begin{proof}
L'identité \(J_W^2=-I\) vérifie les règles de multiplication de
\(\mathbb F_9\). Comme
\([\mathbb F_9:\mathbb F_3]=2\) et
\(\dim_{\mathbb F_3}V_W=6\), la dimension sur \(\mathbb F_9\) est trois.
\end{proof}

Le même opérateur organise le code ternaire au moyen du graphe
\begin{equation}
 C_W=\{(q,qA_W):q\in\mathbb F_3^6\}\subset\mathbb F_3^{12}.
 \label{eq:codigo-grafo-J-anexo}
\end{equation}
Si \(G=(I_6\mid A_W)\), alors
\[
 GG^{\mathsf T}=I_6+A_W A_W^{\mathsf T}=0.
\]
Ainsi, la racine interne de moins l'identité participe simultanément à la
structure quadratique et à l'isotropie du code. Cette double fonction explique
pourquoi la phase et l'incidence exceptionnelle partagent un antécédent et
conservent des lecteurs différents.

\section{Réalification orthogonale de la phase}

La phase analytique se construit sur un espace réel avant le choix d'une carte
complexe. Soit \((E,g)\) un espace euclidien réel de dimension paire et soit
\(J_E\in\operatorname{End}(E)\) un opérateur tel que
\begin{equation}
 J_E^2=-I,
 \qquad
 g(J_Eu,J_Ev)=g(u,v).
 \label{eq:estructura-compleja-ortogonal-anexo}
\end{equation}
La deuxième identité implique \(J_E^{\mathsf T}=-J_E\).

\begin{proposition}[Phase comme groupe orthogonal à un paramètre]
\label{prop:fase-ortogonal-J-anexo}
Pour tout \(\theta\in\mathbb R\),
\begin{equation}
 R_J(\theta):=e^{\theta J_E}
 =\cos\theta\,I+\sin\theta\,J_E
 \label{eq:rotacion-J-anexo}
\end{equation}
est orthogonal et satisfait
\begin{equation}
 R_J(\theta)R_J(\phi)=R_J(\theta+\phi),
 \qquad R_J(2\pi)=I.
 \label{eq:grupo-fase-J-anexo}
\end{equation}
\end{proposition}

\begin{proof}
La série exponentielle se sépare en puissances paires et impaires et, en utilisant
\(J_E^2=-I\), produit l'expression trigonométrique. L'antisymétrie de
\(J_E\) donne
\(R_J(\theta)^{\mathsf T}=R_J(-\theta)\) ; par conséquent,
\(R_J(\theta)^{\mathsf T}R_J(\theta)=I\). La loi de composition résulte du fait
que tous les facteurs sont des fonctions du même opérateur.
\end{proof}

Le choix de la carte \(J_E\leftrightarrow i\) transforme \(E\) en un
espace vectoriel complexe : pour \(a,b\in\mathbb R\),
\begin{equation}
 (a+ib)u:=au+bJ_Eu.
 \label{eq:carta-compleja-J-anexo}
\end{equation}
Dans cette carte, \(R_J(\theta)\) s'écrit comme multiplication par
\(e^{i\theta}\). La notation complexe publie la phase que l'opérateur réel
transporte déjà.

\begin{remark}[Deux échelles, une même relation quadratique]
L'opérateur \(J_W\) agit sur \(\mathbb F_3^6\) ; l'opérateur \(J_E\) agit
sur un espace réel muni d'un produit scalaire. Le prolongement
\(J_W\rightsquigarrow J_E\) est une réalisation qui conserve la relation
\(J^2=-I\) et déclare en outre topologie, norme et complétion. Cet ordre évite
d'utiliser la structure analytique continue comme prémisse de l'incidence finie.
\end{remark}

\section{La droite projective complexe et la sphère de Bloch}

Soit \(\mathcal H_J\) un espace de Hilbert de dimension complexe deux obtenu
au moyen de la carte \eqref{eq:carta-compleja-J-anexo}. Un état pur normalisé
est un vecteur
\begin{equation}
 \psi=\begin{pmatrix}z_0\\ z_1\end{pmatrix},
 \qquad |z_0|^2+|z_1|^2=1,
 \label{eq:espinor-normalizado-anexo}
\end{equation}
et son état physique projectif est le rayon
\([\psi]=\{e^{i\vartheta}\psi:\vartheta\in\mathbb R\}\).
L'espace des rayons est \(\mathbb P(\mathcal H_J)\simeq\mathbb{CP}^{1}\).

Nous définissons les coordonnées
\begin{equation}
 \begin{aligned}
 r_1&=2\operatorname{Re}(\overline z_0z_1),\\
 r_2&=2\operatorname{Im}(\overline z_0z_1),\\
 r_3&=|z_0|^2-|z_1|^2.
 \end{aligned}
 \label{eq:coordenadas-bloch-anexo}
\end{equation}

La représentation sphérique des états à deux niveaux se formule dans les
coordonnées introduites par Bloch \cite{Bloch1946}.

\begin{theorem}[Identification projective de Bloch]
\label{thm:CP1-esfera-bloch-anexo}
L'application
\begin{equation}
 \mathcal B:\mathbb{CP}^{1}\longrightarrow S^2,
 \qquad [\psi]\longmapsto(r_1,r_2,r_3),
 \label{eq:mapa-bloch-anexo}
\end{equation}
est une bijection lisse. Le projecteur de rang un associé à \([\psi]\) est
\begin{equation}
 \rho_\psi=|\psi\rangle\langle\psi|
 =\frac12\bigl(I_2+r_1\sigma_1+r_2\sigma_2+r_3\sigma_3\bigr),
 \label{eq:proyector-bloch-anexo}
\end{equation}
où \(\sigma_1,\sigma_2,\sigma_3\) sont les matrices de Pauli.
\end{theorem}

\begin{proof}
La multiplication de \(\psi\) par une phase commune laisse invariantes les trois
coordonnées. De plus,
\[
 r_1^2+r_2^2+r_3^2
 =4|z_0|^2|z_1|^2+(|z_0|^2-|z_1|^2)^2=1.
\]
Tout \(\boldsymbol r\in S^2\) détermine le projecteur positif
\(\rho=(I+\boldsymbol r\cdot\boldsymbol\sigma)/2\), qui satisfait
\(\rho^2=\rho\) et \(\operatorname{tr}\rho=1\) ; son image est un unique rayon.
La formule matricielle directe de \(|\psi\rangle\langle\psi|\) fournit
\eqref{eq:proyector-bloch-anexo}. Les cartes affines usuelles de
\(\mathbb{CP}^{1}\) rendent l'application et son inverse lisses.
\end{proof}

La sphère de Bloch enregistre la classe projective et conserve, à travers
\(\rho_\psi\), toutes les espérances d'observables hermitiens à deux niveaux.
La phase globale appartient à la fibre \(S^1\) du fibré de Hopf
\cite{Hopf1931}
\begin{equation}
 S^1\longrightarrow S^3\longrightarrow S^2.
 \label{eq:hopf-bloch-anexo}
\end{equation}
La distinction HMT entre état enrichi et coordonnée publiée trouve
ici une réalisation précise : le vecteur normalisé appartient à \(S^3\), le
lecteur projectif publie un point de \(S^2\), et la fibre conserve la phase que
la projection identifie.

\section{Action spinorielle, rotations et transformations de Möbius}

Tout élément de \(\mathrm{SU}(2)\) peut s'écrire comme
\begin{equation}
 U=
 \begin{pmatrix}
 \alpha&\beta\\
 -\overline\beta&\overline\alpha
 \end{pmatrix},
 \qquad |\alpha|^2+|\beta|^2=1.
 \label{eq:SU2-parametrizacion-anexo}
\end{equation}
Il agit linéairement sur les spineurs et, par conjugaison, sur les projecteurs.

\begin{theorem}[Revêtement double et action sur la sphère]
\label{thm:SU2-SO3-bloch-anexo}
Pour chaque \(U\in\mathrm{SU}(2)\), il existe une unique
\(R_U\in\mathrm{SO}(3)\) telle que
\begin{equation}
 U(\boldsymbol r\cdot\boldsymbol\sigma)U^\dagger
 =(R_U\boldsymbol r)\cdot\boldsymbol\sigma.
 \label{eq:accion-SU2-bloch-anexo}
\end{equation}
L'application \(U\mapsto R_U\) est un homomorphisme surjectif de noyau
\(\{\pm I_2\}\).
\end{theorem}

\begin{proof}
La conjugaison conserve l'espace réel tridimensionnel des matrices hermitiennes
sans trace et leur produit scalaire
\(\langle A,B\rangle=\tfrac12\operatorname{tr}(AB)\). Elle induit donc une
transformation orthogonale de \(\mathbb R^3\). La continuité et la valeur en
l'identité fixent un déterminant positif. Si la conjugaison est triviale sur les
trois matrices de Pauli, \(U\) commute avec toute matrice complexe \(2\times2\),
et, par le lemme de Schur, est scalaire ; la condition de déterminant un laisse
\(U=\pm I_2\). La représentation axe--angle montre que toute rotation possède
deux antécédents opposés.
\end{proof}

Dans la carte affine \(z=z_0/z_1\) de \(\mathbb{CP}^{1}\), la même action prend
la forme
\begin{equation}
 z\longmapsto
 \frac{\alpha z+\beta}
      {-\overline\beta z+\overline\alpha}.
 \label{eq:mobius-SU2-anexo}
\end{equation}
La rotation de Bloch et la transformation fractionnaire sont donc deux
coordonnées de la même action projective. La première emploie le projecteur
hermitien ; la deuxième, la sphère de Riemann.

\section{La spécialisation icosaédrique et le groupe binaire \texorpdfstring{\(2A_5\)}{2A5}}

Le groupe rotationnel de l'icosaèdre est un sous-groupe de
\(\mathrm{SO}(3)\) isomorphe à \(A_5\). Son image réciproque par le revêtement du
\cref{thm:SU2-SO3-bloch-anexo} définit le groupe icosaédrique binaire
\begin{equation}
 2A_5:=\pi^{-1}(A_5)\subset\mathrm{SU}(2).
 \label{eq:def-2A5-anexo}
\end{equation}

\begin{proposition}[Extension centrale icosaédrique]
\label{prop:extension-central-2A5-anexo}
Il existe une suite exacte
\begin{equation}
 1\longrightarrow\{\pm I_2\}
 \longrightarrow2A_5
 \longrightarrow A_5
 \longrightarrow1,
 \label{eq:sucesion-2A5-anexo}
\end{equation}
et \(|2A_5|=120\). Une rotation icosaédrique d'ordre trois ou cinq se relève en
des éléments d'ordre six ou dix, respectivement.
\end{proposition}

\begin{proof}
La restriction de \(\pi:\mathrm{SU}(2)\to\mathrm{SO}(3)\) à l'image réciproque de
\(A_5\) est surjective et conserve le noyau \(\{\pm I_2\}\). L'exactitude et
le théorème de l'ordre donnent \(|2A_5|=2|A_5|=120\). Une rotation d'angle
\(2\pi/n\) est représentée dans \(\mathrm{SU}(2)\) avec un demi-angle
\(\pi/n\) ; après \(n\) itérations, elle atteint \(-I_2\) et, après
\(2n\), l'identité.
\end{proof}

La spécialisation icosaédrique fixe une articulation rigoureuse entre la géométrie des
états à deux niveaux, les transformations de Möbius et la branche des symétries
finies. Le développement de McKay et des structures exceptionnelles utilise
ensuite la représentation de \(2A_5\) et ses modules ; son contenu supplémentaire
est développé dans les théorèmes spécifiques de cette branche.

\section{Composition fonctorielle \texorpdfstring{\(\mathbb F_3^6\to\mathbb F_9^3\to\mathbb{CP}^{1}\simeq S^2\to\mathrm{SU}(2)\to\mathrm{SO}(3)\to 2A_5\)}{F3^6 -> F9^3 -> CP1 -> S2 -> SU(2) -> SO(3) -> 2A5}}

La relation \(J^2=-I\) traverse quatre catégories en conservant son type :
\begin{enumerate}
 \item dans \(\mathbb F_3^6\), elle définit l'extension quadratique
       \(\mathbb F_9\) et organise l'incidence de Witt ;
 \item dans un espace réel orthogonal, elle engendre le flux de phase
       \(e^{\theta J_E}\) ;
 \item en dimension complexe deux, la projectivisation produit
       \(\mathbb{CP}^{1}\simeq S^2\) et les projecteurs de Bloch ;
 \item l'action de \(\mathrm{SU}(2)\) descend aux rotations et aux
       transformations de Möbius, et son image réciproque icosaédrique produit
       \(2A_5\).
\end{enumerate}
La structure complexe interne possède ainsi une provenance finie et une
réalisation analytique explicite. La sphère de Bloch publie la géométrie des
rayons ; la généalogie HMT conserve en outre le chemin, l'orientation, la retenue et la
mémoire qui précèdent cette publication.
